\section{Introduction}\label{sec:introduction}
Investor opinions and managerial disclosures can both contain information relevant to stock prices. Studies of investor text examine expressed beliefs on financial social media \citep{chen2014,bartov2018,cookson2024}, while studies of corporate language examine the information and beliefs conveyed by managers \citep{li2010,jiang2019}. Considering either side alone leaves a question unanswered: does the combination of an investor's concern and the corresponding managerial reply convey return information beyond the two tones separately? A positive reply to an anxious, specific question need not have the same meaning as equally positive language in a generic exchange.

Interactive platforms make this question observable. Investors choose what to ask, and managers decide whether and how to respond. Inquiries can elicit additional details, while the specificity of replies can help readers interpret management's language \citep{hollander2010,matsumoto2011}. At the same time, questions and replies can react to previous prices or reflect strategic disclosure. We therefore distinguish the predictive association of a dialogue measure from evidence that replies cause price changes or disclose previously unknown fundamentals.

The Chinese market provides a setting in which to compare information in public dialogue with possible trading pressure. The Shenzhen and Shanghai exchanges operate interaction platforms covering a broad cross section of listed firms \citep{lee2022shall}. Substantial retail participation connects these communications to investor trading, and short-sale constraints may limit the offsetting response of arbitrageurs \citep{liao2022,jiang2022}. Neither observation by itself establishes irrational sentiment. Chinese prices can also contain information about future fundamentals \citep{carpenter2021}. The distinction between information and pressure is consequently an empirical question rather than an assumption about the market.

We construct the Effective Sentiment Index (ESI) from matched investor inquiries and managerial replies. Investor tone supplies the signed investor component, scaled by question volume. The managerial component combines standardized reply tone with reply substantiveness and a smooth weight determined by the preceding investor tone. The same managerial tone can therefore receive different weights depending on what investors asked and how specifically managers responded. The weights are maintained modeling choices; their economic motivation does not establish that they are optimal.

The historical question-date analysis reproduces the Annual Review program using its frozen rules. Equal-weighted (EW) and value-weighted (VW) ESI portfolios earn gross daily O2C excess returns of 5.31 and 5.26 basis points, respectively. Under the original arithmetic annualization convention these are 12.85\% and 12.72\%. Both exceed the corresponding point estimates for investor tone or manager tone alone. ESI also has a positive one-day regression coefficient after the original controls and fixed effects. These comparisons describe historical associations: the original pools and factor alignment are retained, and replies recorded after the question date enter that date's signal.

We separately examine the information clock using existing public-availability C2C results. In a common sample, the standalone ESI coefficient is 1.48 basis points ($t=4.36$), but falls to 0.58 ($t=0.75$) after including the complete set of text features. On common portfolio dates and a common candidate pool, the EW China four-factor alpha is 4.68 basis points ($t=3.29$), while the VW alpha is 3.78 ($t=1.54$). These results do not support a stable advantage over manager tone. The C2C extension uses a different frozen calibration and later sample, so the change cannot be attributed solely to reply timing.

Two additional analyses qualify the interpretation. Signals observed before actual earnings announcements do not significantly predict standardized unexpected earnings or announcement returns in the available tests. Full-corpus automated classification describes 19 content themes and five question natures, but joint tests do not reject equal ESI slopes across topics. A significant products-or-business slope in the baseline topic regression loses significance after both tones enter. We retain these results because they distinguish the positive historical association from the stronger claims needed to identify a mechanism.

The study contributes in three ways. First, it develops a dialogue-based sentiment measure that makes attention, reply specificity and the investor state explicit. This differs from treating a reply as standalone managerial language or simply summing two tone scores. Second, it compares the measure with its component signals and separates retrospective O2C evidence from public-availability regressions and common-pool portfolios. This clarifies which conclusions depend on the information clock and competing text information. Third, it brings content and question-nature classifications into the interpretation of interactive disclosures. These classifications describe what investors ask and what managers discuss; they are not taken as independently validated semantic labels or evidence of different priced risks.

Section~\ref{sec:literature_review} reviews related work. Section~\ref{sec:data} describes the data and ESI. Section~\ref{sec:baseline_return_predictability} reports historical return evidence. Section~\ref{sec:supplements} presents the public-availability, earnings and topic evidence, and Section~\ref{sec:conclusion} concludes.
